\section{\systemname: Behavioral Patch Review for Skills}
\label{sec:system}

\systemname{} supports skill repair as a review of three related changes. The
\emph{intended behavioral delta} describes what a workflow owner wants to
change, preserve, or leave unresolved. The \emph{artifact delta} is the
revision encoded in the skill. The \emph{enacted behavioral delta} is the
change observed when the original and revised skills are executed in
comparable situations. A repair becomes reviewable when these three deltas can
be examined together rather than inferred from a textual edit or a successful
rerun of the motivating incident alone.

Figure~\ref{fig:system-overview} summarizes the interaction. The owner first
constructs a behavioral scope from concrete cases. The system then organizes
evidence that relates this scope to the current skill and to a candidate
revision. Matched executions make the behavioral reach of the candidate
visible before the owner decides whether to release it. These activities are
iterative: evidence may prompt a revision to the intended scope, the skill, or
the cases considered relevant. The system therefore provides a shared review
record rather than treating repair as a one-shot generation task.

\begin{figure*}[t]
    \centering
    \small
    \setlength{\tabcolsep}{4pt}
    \begin{tabular}{@{}C{0.225\textwidth} C{0.225\textwidth} C{0.225\textwidth} C{0.225\textwidth}@{}}
        \fbox{\begin{minipage}[t][3.55cm][t]{0.205\textwidth}
            \textbf{1. Articulate intent}\\[3pt]
            Examine the failure and nearby cases\\
            Record behavior to change, preserve,\\
            or leave unresolved\\[4pt]
            Construct the intended behavioral delta.
        \end{minipage}}
        &
        \fbox{\begin{minipage}[t][3.55cm][t]{0.205\textwidth}
            \textbf{2. Assemble evidence}\\[3pt]
            Relate commitments to source regions\\
            Inspect relevant executions\\
            Review evidentiary limits\\[4pt]
            Connect intent to the current skill.
        \end{minipage}}
        &
        \fbox{\begin{minipage}[t][3.55cm][t]{0.205\textwidth}
            \textbf{3. Inspect the revision}\\[3pt]
            Generate, edit, or import a candidate\\
            Review the textual delta with links\\
            to intent and evidence\\[4pt]
            Make the artifact delta legible.
        \end{minipage}}
        &
        \fbox{\begin{minipage}[t][3.55cm][t]{0.205\textwidth}
            \textbf{4. Review enacted change}\\[3pt]
            Compare original and revised behavior\\
            Identify intended effects, regressions,\\
            and unresolved boundaries\\[4pt]
            Release, revise, investigate, or defer.
        \end{minipage}}
    \end{tabular}
    \caption{\systemname{} makes the intended, artifact, and enacted deltas
    jointly reviewable. Evidence connects the owner's behavioral scope to the
    current skill and candidate revision; matched executions support revision
    and release decisions.}
    \Description{Four panels show articulating intended behavior, assembling
    execution evidence, inspecting a skill revision, and reviewing enacted
    behavioral change before a release decision.}
    \label{fig:system-overview}
\end{figure*}

\subsection{Interaction Overview}
\label{sec:system-walkthrough}

We illustrate the workflow through Lin, the researcher introduced in
Section~\ref{sec:introduction}. Her travel-recovery skill selected an
inexpensive replacement flight that arrived after a fixed conference
presentation. The execution was plainly unacceptable, but it did not by
itself establish the policy that should govern future disruptions.

Lin begins by expressing the desired change in her own terms: fixed
commitments should take priority over price. \systemname{} places this incident
alongside neighboring cases that expose the reach of such a rule. Lin wants to
retain cost sensitivity for routine travel and confirmation before a
consequential purchase, while leaving an extreme fare trade-off unresolved.
Together, these judgments turn an objection to one outcome into a provisional
behavioral scope.

The system next brings together relevant parts of the current skill and
evidence from controlled executions. A targeted source intervention suggests
that an instruction treating calendar conflicts as a soft trade-off is
behaviorally relevant in the incident. The result is presented as bounded
evidence about sensitivity to a source region, not as proof of unique causal
responsibility or validation of a future revision. Lin can revise her scope if
this evidence changes her interpretation of the problem.

An AI service then drafts a candidate from the current skill, Lin's recorded
scope, and the reviewed evidence. \systemname{} presents the candidate as a
line-based diff. Annotations connect each changed block to the behavioral
commitment it is intended to address and to the evidence that informed the
revision. These links make the rationale and provenance of the artifact delta
inspectable without implying that a textual change has achieved its intended
effect.

Finally, \systemname{} executes the complete original and revised skills on
the reviewed cases under matched conditions. Lin can see whether the incident
is repaired, whether valued behaviors are preserved, and whether the revision
enters a boundary for which she has not yet defined a policy. She may revise
the skill or behavioral scope, request further evidence, defer the repair, or
release the reviewed version. The same scope therefore links her initial
intent to the textual revision, enacted outcomes, and final decision.

\subsection{Defining the Behavioral Scope of a Skill Repair}
\label{sec:behavioral-scope}

A review begins with an observable breakdown, the active skill, and the
available execution record. \systemname{} brings the task context, skill
instructions, tool activity, decisions, and outcome into a common view. This
reconstruction helps the owner distinguish what the skill said from what the
agent enacted and provides a stable point of comparison for later evidence.

The owner constructs a behavioral scope from the motivating incident,
accepted prior behavior, and contrastive cases. Each case $c$ combines a task
context, baseline evidence $B_0(c)$, and one of three normative dispositions:

\begin{itemize}[leftmargin=*]
    \item \textsc{Change}$(g_c)$ identifies behavior the revision should
    produce, such as selecting an itinerary that arrives before a fixed
    presentation.
    \item \textsc{Preserve}$(p_c)$ identifies a valued property that should
    remain stable, such as requesting confirmation before a consequential
    purchase.
    \item \textsc{Unresolved} marks a policy boundary for which the owner is
    not yet prepared to authorize a general rule, such as the acceptable fare
    premium for protecting a commitment.
\end{itemize}

Cases that are relevant to monitoring but outside the current repair can be
kept outside the candidate author's normative input. If a revision changes one
of these cases, the system returns it to the owner for judgment rather than
treating it as evidence of success or failure. This distinction separates
uncertainty about policy from evidence that an authorized commitment was met.

Commitments can concern decisions, tool calls, action ordering, confirmation
points, or outcomes. Each is paired with an observable criterion that the
owner can inspect. The system preserves successive versions of the scope so
that later evidence can be interpreted against the commitments that were in
place when a candidate was produced.

\paragraph{Scope anchors and contrastive cases.}
The motivating incident anchors the repair, but nearby cases clarify its
intended reach. These cases can come from prior executions, examples supplied
by the owner, or system-generated contrasts reviewed by the owner. Each
contrast makes explicit how it differs from an existing case and which policy
boundary it probes. This comparative representation allows a vague request
such as ``prioritize important events'' to become a set of concrete, revisable
commitments without requiring the owner to author a complete policy upfront.

\paragraph{Pre-reveal boundary commitments.}
Where possible, the owner records these commitments before seeing candidate
outcomes. This ordering preserves a prospective account of intended behavior
rather than allowing the observed candidate to redefine the repair
retrospectively. Selected cases may also remain outside an AI author's input,
allowing later executions to examine whether the revision extends beyond the
examples that directly informed it. These cases do not certify generalization;
they provide bounded evidence about the reach of the reviewed repair.

\subsection{Locating Source Instructions and Reviewing the Repair Preview}
\label{sec:instruction-effects}

To help the owner relate an observed breakdown to the skill, \systemname{}
supports targeted interventions on an isolated copy of the source artifact. A
model-assisted planner identifies an instruction and an already reconstructed
case that are relevant to the current scope; the execution service then
compares the original behavior with behavior under bounded variants of that
instruction. The reviewed skill itself is not changed during this process.

The resulting evidence answers a deliberately narrow question: under the
selected case and execution conditions, is observable behavior sensitive to
this source region? It does not establish unique causality, capture
interactions among all instructions, or validate text that has not yet been
executed. \systemname{} keeps this limitation visible and retains links to the
underlying artifact versions, cases, executions, and observable differences.

The repair preview places this source-location evidence alongside the owner's
behavioral scope. Its purpose is to let the owner review the evidence and
commitments that will inform implementation, not to ask them to approve a
system-generated policy. The owner can return to the scope whenever the
evidence reveals that the intended repair was underspecified. Exact
interventions and execution budgets for the study instantiation are reported
in Appendix~\ref{app:probe-protocol}.

\subsection{Reviewing the Skill Diff and Its Behavioral Consequences}
\label{sec:candidate-review}

The workflow owner may write the candidate skill or delegate its
implementation to an AI service or collaborator. Regardless of authorship,
\systemname{} records the candidate as a revision of the current skill and
associates it with the behavioral scope and evidence available at the time.
This provenance preserves the distinction between deciding what behavior is
acceptable and choosing text intended to produce it.

The primary artifact view uses a familiar line-based review layout. Added and
removed skill text is accompanied by compact annotations that distinguish four
kinds of information:

\begin{itemize}[leftmargin=*]
    \item \textbf{Scope commitment} identifies the intended behavior that a
    changed block is meant to address.
    \item \textbf{Implementation rationale} records why the candidate author
    associated that text with the commitment.
    \item \textbf{Source evidence} shows execution evidence that made the
    corresponding region of the original skill relevant to the repair.
    \item \textbf{Candidate evidence}, when available, reports bounded
    evidence about the behavioral sensitivity of the revised block.
\end{itemize}

The annotations support navigation between the behavioral and textual objects
of review; they do not attribute a whole-candidate outcome to a single edited
block. Evidence about the source skill, an optional intervention on candidate
text, and execution of the complete candidate therefore remain visibly
distinct.

After a candidate is committed, \systemname{} executes the complete original
and revised skills on the reviewed cases. For case $c$, it estimates candidate
behavior $B_p(c)$ under conditions matched to the baseline $B_0(c)$, including
task context, tool state, model configuration, and side-effect policy. The
comparison evaluates the candidate as a whole and organizes the result into
four review states:

\begin{table}[t]
    \caption{Primary review states and their meanings.}
    \label{tab:review-states}
    \small
    \begin{tabularx}{\columnwidth}{@{}p{0.25\columnwidth}Y@{}}
        \toprule
        \textbf{State} & \textbf{Meaning} \\
        \midrule
        Aligned & Available evidence is consistent with a recorded Change or Preserve commitment. \\
        Mismatch & A desired change remains unmet or a valued behavior regresses. \\
        Needs judgment & The candidate changes behavior in a region for which the owner has not authorized a target. \\
        Insufficient evidence & The available executions do not support a stable comparison. \\
        \bottomrule
    \end{tabularx}
\end{table}

The workspace gives priority to mismatches, unresolved changes, and
insufficient evidence while keeping commitments separate. Repairing the
motivating incident therefore cannot conceal a regression in a behavior the
owner intended to preserve. Cross-highlighting connects a selected case to its
commitment, whole-candidate outcome, and related skill blocks, helping the
owner move between intended, artifact, and enacted deltas.

When a result remains difficult to interpret, the owner can inspect individual
executions, gather additional evidence, or request a bounded intervention on a
candidate block. These actions provide further evidence without converting an
unresolved policy question into an automated judgment. The owner can then edit
the candidate, revise the scope, defer the repair, or release the reviewed
version. Cases and criteria accepted during review can inform future
regression checks.

\paragraph{Behavior evaluation.}
The evaluated prototype represents behavior through observable tool traces and
environment outcomes. Fact criteria capture properties such as arrival time or
ledger creation; trace criteria capture events and ordering, such as
confirmation preceding booking. In the controlled study, confirmatory outcomes
use deterministic rules over these reconstructable records rather than an LLM
semantic judge. Cases without an authorized target are evaluated only for
material behavioral change and returned to the owner for judgment. The system
stores observable traces and outcomes, not hidden chain-of-thought.

\subsection{System Architecture, Implementation, and Scope}
\label{sec:system-implementation}

\systemname{} is implemented as a web-based review workspace backed by an
orchestration service. A skill comprises natural-language instructions, a
typed tool interface, and an execution configuration. Each execution receives
an isolated copy of a stateful tool environment, allowing the system to expose
realistic decisions, tool results, and side effects without reaching a
production account. Original, intervened, and candidate skills are retained as
separate versions so that evidence can be traced to the artifact and context
that produced it.

LLM services perform bounded interpretive and generative work: they help
retrieve relevant cases, propose source regions for inspection, draft
candidate revisions, and explain computed comparisons. They do not fabricate
tool results, modify recorded owner commitments, determine confirmatory
outcomes, or authorize release. Execution and evaluation services maintain
these boundaries by validating tool calls against the scenario environment and
computing study outcomes from observable traces and end states.

The study instantiation freezes model configurations, scenario worlds,
candidate-generation inputs, and execution budgets within each condition.
It also records the provenance of scope revisions, candidate versions,
executions, and release decisions. These controls support matched comparisons
and reconstruction of a review without turning the product interaction into an
experiment-specific sequence. Appendix~\ref{app:system-details} reports the
data model, intervention protocol, execution configuration, and prompt
contracts used in the evaluated prototype.

\systemname{} provides bounded evidence about a revision rather than a
guarantee of behavior outside the reviewed cases and execution conditions.
Targeted interventions indicate sensitivity, not unique causal responsibility;
matched executions characterize the complete candidate only within the
observed contexts. The workflow therefore retains explicit states for
unresolved policy and insufficient evidence, and allows the owner to keep the
current skill whenever a revision is not yet justified.
